\documentclass[aps,pre,twocolumn,showpacs,amsmath,amssymb]{revtex4-1}

\usepackage{graphicx}
\usepackage{color}

\usepackage[portuguese]{babel}
\usepackage[utf8]{inputenc}
\usepackage[T1]{fontenc}

\hfuzz 1pt
\vfuzz 1pt

\setlength{\parskip}{\baselineskip}

\begin{document}

\title{Exercício 00: inserir nome e número do exercício}

\author{inserir nome e número do aluno}
\email{inserir email}

\begin{abstract} 
  Sumário do trabalho: máximo de 3 linhas. Descrição muito breve dos objetivos do exercício.
\end{abstract}

\maketitle

Relatório do trabalho:

\begin{itemize}
%
\item descrever os pârametros utilizados: o relatório deve ser auto-consistente para que qualquer pessoa possa reproduzir os resultados sem necessitar de informação extra;
\item figuras com resultados podem ser incluídas com uma legenda completa que permita perceber a figura sem recorrer ao texto principal.
\item o relatório está limitado a 2 páginas neste formato (revtex e duas colunas);
\item o relatório deve ser enviado em pdf para o email do docente até à data limite (ver folha de exercício);
\item em anexo ao relatório deve ser enviado o código (.cpp) sem qualquer outro ficheiro adicional.
%
\end{itemize}

Deve ser incluída uma página extra com a declaração do uso de Inteligência Artificial,  descrevendo detalhadamente como foi utilizada.

\end{document}
